
\medskip
\phantomsection\label{rem:sga3-VIB-11.18.5}%
\noindent\textbf{Remarks 11.18.5.---}
\begin{enumerate}
\item[\textup{(a)}]\phantomsection\label{rem:sga3-VIB-11.18.5-a}%
  One consequence of the preceding theorem is that the category
  of commutative affine \(k\)-groups is abelian. For this, as well
  as for other results on affine \(k\)-groups, we refer to
  \textup{[DG70]}, \S~III.3, \SGARefLink{prop:sga3-VIB-7.4}{7.4}--\SGARefLink{prop:sga3-VIB-7.8}{7.8}.

\item[\textup{(b)}]\phantomsection\label{rem:sga3-VIB-11.18.5-b}%
  Finally, let us point out that M.~Takeuchi gave another proof
  of the results of \hyperref[thm:sga3-VIB-11.17]{11.17}--\hyperref[cor:sga3-VIB-11.18.4]{11.18.4}; see \textup{[Ta72]}, \S~5.
  He moreover strengthened \hyperref[concl:sga3-VIB-11.18.4-i]{11.18.4}\textup{\SGARefLink{cor:sga3-VIB-11.18.4:item:i}{(i)}} above by showing
  that \(A\) is even a projective \(B\)-module; see
  \textup{[Ta79]}, \SGARefLink{sec:sga3-VIB-5}{Theorem~5} \textup{(see also [MW94],
  Theorem~3.6)}.
\end{enumerate}

\section*{12. Complements on \(G_{\mathrm{af}}\) and the
  ``anti-affine'' groups}
\phantomsection\label{sec:sga3-VIB-12}%
\addcontentsline{toc}{section}{12. Complements on Gaf and the
  ``anti-affine'' groups}

\footnote{\label{note:sga3-VIB-135}Editor's note: this
section has been added.}
We begin with the following lemma, which extends
\hyperref[lem:sga3-VIB-11.18.1]{11.18.1} to the case where
\(G\) is not assumed to be of finite type.%

\medskip
\phantomsection\label{lem:sga3-VIB-12.1}%
\noindent\textbf{Lemma 12.1.---}
\emph{Let \(k\) be a field and let
\(u:G\to H\) be a monomorphism of \(k\)-groups, with \(H\)
affine. Suppose that \(u\) is quasi-compact. Then \(u\) is a
closed immersion.}

By \textup{(fpqc)} descent, we may suppose that \(k\) is
algebraically closed. By
\hyperref[prop:sga3-VIA-6.4]{VI\(_A\), 6.4}, the closed image
\(I\) of \(u\) is a closed group subscheme of \(H\), hence is
still affine. Replacing \(H\) by \(I\), we may therefore suppose
that \(u\) is schematically dominant. Since \(k\) is
algebraically closed, \(H'=H_{\mathrm{red}}\) is a group
subscheme of \(H\). Put \(G'=G\times_HH'\). The morphism
\[
  u':G'\longrightarrow H'
\]
deduced from \(u\) by base change is a quasi-compact
monomorphism and is dominant \textup{(the underlying continuous
map being the same for \(u\) and \(u'\))}. Hence, by
\hyperref[cor:sga3-VIA-6.2]{VI\(_A\), 6.2}, \(u'\) is faithfully
flat. It is therefore a quasi-compact faithfully flat
monomorphism, and thus an isomorphism \textup{(cf.\ IV, \SGARefLink{sga3:iv:remark:1:14}{1.14})}.

Consequently \(u:G\to H\) is a homeomorphism, and hence is
affine by \hyperref[lem:sga3-VIB-2.9.1]{2.9.1}. Thus \(G\) is
affine, and \(u\) is a closed immersion by
\hyperref[thm:sga3-VIB-11.18.2]{11.18.2\textup{(iii)}}.

\medskip
\phantomsection\label{thm:sga3-VIB-12.2}%
\noindent\textbf{Theorem 12.2.---}
\emph{Let \(G\) be an algebraic \(k\)-group. Denote by
\(\rho:G\to G_{\mathrm{af}}\) the canonical morphism and by
\(N\) its kernel.}
\begin{enumerate}
\item[\textup{(i)}]\phantomsection\label{thm:sga3-VIB-12.2-i}%
  \emph{The canonical morphism
  \(G/N\to G_{\mathrm{af}}\) is an isomorphism. Consequently
  \(G_{\mathrm{af}}\) is an affine algebraic group and
  \(\rho\) is faithfully flat.}

\item[\textup{(ii)}]\phantomsection\label{thm:sga3-VIB-12.2-ii}%
  \emph{There is a canonical isomorphism
  \((G/N)_{\mathrm{af}}=G_{\mathrm{af}}\).}

\item[\textup{(iii)}]\phantomsection\label{thm:sga3-VIB-12.2-iii}%
  \emph{\(N\) is a characteristic subgroup of \(G\).}

\item[\textup{(iv)}]\phantomsection\label{thm:sga3-VIB-12.2-iv}%
  \emph{\(\mathcal O(N)=k\).}

\item[\textup{(v)}]\phantomsection\label{thm:sga3-VIB-12.2-v}%
  \emph{\(N\) is smooth, connected, and commutative.}
\end{enumerate}

Assertion~\hyperref[thm:sga3-VIB-12.2-i]{\textup{(i)}} is a
special case of
\hyperref[lem:sga3-VIB-11.18.1]{11.18.1}, and
\hyperref[thm:sga3-VIB-12.2-ii]{\textup{(ii)}} follows from the
universal properties of \(G_{\mathrm{af}}\) and
\((G/N)_{\mathrm{af}}\).

Let us prove
\hyperref[thm:sga3-VIB-12.2-iii]{\textup{(iii)}}. For an
arbitrary \(k\)-scheme
\(\pi:S\to\operatorname{Spec}k\), consider the cartesian square
\begin{center}
  \phantomsection\label{diag:sga3-VIB-12.2-base-change}%
  \begin{tikzcd}[row sep=large,column sep=large]
    G_S \arrow[r] \arrow[d,"q"'] &
    G \arrow[d,"p"] \\
    S \arrow[r,"\pi"] &
    \operatorname{Spec}k\mathpunct{,}
  \end{tikzcd}
\end{center}
Since \(p\) is quasi-compact and separated and \(\pi\) is flat,
EGA~III, 1.4.15 and EGA~IV\(_1\), 1.7.21 give
\phantomsection\label{form:sga3-VIB-12.2-base-change-global-sections}%
\[
  q_*(\mathcal O_{G_S})
   =\pi^*(\mathcal O(G))
   =\pi^*(\mathcal O(G_{\mathrm{af}})).
\]
Hence EGA~II, 1.5.2 gives
\[
  (G_S)_{\mathrm{af}}=(G_{\mathrm{af}})_S.
\]
Thus \(N_S\), being the kernel of the canonical morphism
\(G_S\to(G_S)_{\mathrm{af}}\), is invariant under every
automorphism of \(G_S\); in other words, \(N\) is a
characteristic subgroup of \(G\).

To prove \hyperref[thm:sga3-VIB-12.2-iv]{\textup{(iv)}}, put
\[
  N'=\operatorname{Ker}(N\longrightarrow N_{\mathrm{af}}).
\]
By \hyperref[thm:sga3-VIB-12.2-ii]{\textup{(ii)}}, this is an
invariant subgroup of \(G\). Since \(N\) is algebraic
\textup{(being a closed subgroup of \(G\))},
\hyperref[thm:sga3-VIB-12.2-i]{\textup{(i)}} gives
\(N/N'\simeq N_{\mathrm{af}}\). Moreover, by VI\(_A\), \hyperref[thm:sga3-VIA-3.2]{3.2}
and \SGARefLink{par:sga3-VIA-5.3.2}{5.3.2}, there is an isomorphism of \(k\)-groups
\phantomsection\label{form:sga3-VIB-12.2-double-quotient}%
\[
  (G/N')/N_{\mathrm{af}}
    \simeq (G/N')/(N/N')
    \simeq G/N.
\]
Since \(N_{\mathrm{af}}\) is affine, the projection
\(G/N'\to G/N\) is affine by
\hyperref[prop:sga3-VIB-9.2]{9.2\textup{(vii)}}; and since
\(G/N=G_{\mathrm{af}}\) is affine, \(G/N'\) is affine as well.
The universal property of \(G_{\mathrm{af}}\) therefore makes
the projection
\[
  p':G\longrightarrow G/N'
\]
factor through \(G_{\mathrm{af}}=G/N\). Hence \(N\subset N'\),
and consequently \(N=N'\). Thus \(N_{\mathrm{af}}\) is the
trivial group, and \(\mathcal O(N)=k\).

Finally,
\hyperref[thm:sga3-VIB-12.2-v]{\textup{(v)}} follows from the
next lemma.

\medskip
\phantomsection\label{lem:sga3-VIB-12.3}%
\noindent\textbf{Lemma 12.3.---}
\emph{Let \(k\) be a field and let \(N\) be an algebraic
\(k\)-group such that \(\mathcal O(N)=k\). Then \(N\) is
smooth, connected, and commutative.}

Indeed, we may suppose that \(k\) is algebraically closed.
Then
\[
  H=N_{\mathrm{red}}^0
\]
is a \(k\)-subgroup of \(N\), and the quotient \(k\)-scheme
\(X=N/H\) is finite \textup{(hence affine)} over \(k\), by
\hyperref[prop:sga3-VIA-5.5.1]{VI\(_A\), 5.5.1} and
\hyperref[prop:sga3-VIA-5.6.1]{5.6.1}. Since
\(p:N\to X\) is faithfully flat,
\[
  \mathcal O(X)\subset\mathcal O(N)=k.
\]
It follows that \(N=N_{\mathrm{red}}^0\), and hence that \(N\)
is smooth \textup{(VI\(_A\), \hyperref[prop:sga3-VIA-1.3.1]{1.3.1})} and connected.

Let \(Z\) be the center of \(N\). By
\hyperref[cor:sga3-VIB-6.2.6]{6.2.6}, \(N/Z\) is affine, and
the preceding argument gives \(\mathcal O(N/Z)=k\). Hence
\(N=Z\). This proves the lemma and completes the proof of
\hyperref[thm:sga3-VIB-12.2]{12.2}.

We also record, without proof, the following theorem. Recall
that an \emph{abelian variety} over a field \(k\) is a proper,
smooth, connected \(k\)-group scheme.

\medskip
\phantomsection\label{thm:sga3-VIB-12.4}%
\noindent\textbf{Theorem 12.4 \textup{(Chevalley)}.---}
\emph{Let \(k\) be a perfect field and let \(G\) be a smooth,
connected algebraic \(k\)-group. Then there exists a smooth,
connected affine \(k\)-subgroup \(L\), invariant in \(G\),
such that \(G/L\) is an abelian variety. Moreover, \(L\) is
unique and its formation commutes with extension of the base
field.}

\medskip
\phantomsection\label{rem:sga3-VIB-12.5}%
\noindent\textbf{Remarks 12.5.---}
\begin{enumerate}
\item[\textup{(1)}]\phantomsection\label{rem:sga3-VIB-12.5-1}%
  This theorem was announced in 1953 by C.~Chevalley, who
  published his proof in 1960 \textup{([Ch60])}. Meanwhile,
  other proofs were obtained independently by I.~Barsotti and
  M.~Rosenlicht \textup{([Ba55, Ro56])}; see \textup{[Se99]}
  for historical comments.

\item[\textup{(2)}]\phantomsection\label{rem:sga3-VIB-12.5-2}%
  A modern version \textup{(that is, in the language of
  schemes)} of Chevalley's proof was given by B.~Conrad
  \textup{([Co02])}. \textup{(In \emph{loc.\ cit.},
  ``algebraic group'' means a smooth connected \(k\)-group
  scheme.)}

\item[\textup{(3)}]\phantomsection\label{rem:sga3-VIB-12.5-3}%
  A modern version of Rosenlicht's proof, on the other hand,
  was given by Ngô B.-C.\ in a course at Orsay in 2005--2006.

\item[\textup{(4)}]\phantomsection\label{rem:sga3-VIB-12.5-4}%
  If the hypothesis that \(k\) be perfect is omitted, there
  still exists a smallest invariant connected affine subgroup
  \(L\), not necessarily smooth, such that \(G/L\) is an
  abelian variety \textup{([BLR], \S~\SGARefLink{prop:sga3-VIB-9.2}{9.2}, Theorem~1)}.

\item[\textup{(5)}]\phantomsection\label{rem:sga3-VIB-12.5-5}%
  One may also omit the hypothesis that \(G\) be smooth over
  \(k\). Indeed, by VII\(_A\), \SGARefLink{prop:sga3-VIB-8.3}{8.3}, there is an integer
  \(n\geq1\) such that the quotient
  \[
    G'=G/(\operatorname{Fr}^nG)
  \]
  is smooth. Then \(G'\) contains a subgroup \(L'\) as in
  \hyperref[rem:sga3-VIB-12.5-4]{\textup{(4)}} above, and the
  inverse image of \(L'\) in \(G\) again has the same
  properties. Thus, for every connected algebraic group \(G\)
  over a field \(k\), there is an exact sequence
  \phantomsection\label{form:sga3-VIB-12.5-connected-exact-sequence}%
  \[
    1\longrightarrow H\longrightarrow G
      \longrightarrow A\longrightarrow1,
  \]
  where \(H\) is an affine \(k\)-group and \(A\) is an abelian
  \(k\)-variety. Moreover, by \textup{[Per76]},
  Corollary~4.2.9, such an exact sequence exists for every
  connected \(k\)-group \(G\), not necessarily algebraic.

\item[\textup{(6)}]\phantomsection\label{rem:sga3-VIB-12.5-6}%
  Let \(k\) be an algebraically closed field and let \(G\) be
  the semidirect product of an elliptic curve \(E\) by the
  constant \(k\)-group \(\{\pm1\}_k\), for the action defined
  by \((-1)\cdot x=-x\). In this case, if \(L\) is a closed
  invariant subgroup of \(G\) such that \(G/L\) is connected,
  then \(L=G\).
\end{enumerate}

\medskip
\phantomsection\label{rem:sga3-VIB-12.6}%
\noindent\textbf{Remark 12.6.---}
Following \textup{[Br09]}, a \(k\)-group \(N\) will be called
\emph{anti-affine} if \(\mathcal O(N)=k\). By
\hyperref[lem:sga3-VIB-12.3]{12.3} and
\hyperref[thm:sga3-VIB-12.4]{12.4}, if \(k\) is perfect, every
anti-affine algebraic \(k\)-group is an extension of an
abelian variety by a smooth, connected, commutative affine
algebraic \(k\)-group. For the precise structure of
anti-affine algebraic groups over a perfect field, and various
consequences, see the recent articles of M.~Brion and
C.\ \& F.~Sancho de Salas \textup{([Br09, SS09])}.

To conclude this section, we shall prove two results due to
M.~Raynaud, the first being
\hyperref[rem:sga3-VIB-11.11.1]{Remark 11.11.1}, and the
second \SGARefLink{prop:sga3-VIB-2.1}{Proposition~2.1} of Exposé~XVII, Appendix~III. We shall
need the following lemma,%
which improves, for a complete discrete valuation ring \(R\),
the flatness criteria given in \textup{[BAC]}, \S~III.5.

\medskip
\phantomsection\label{lem:sga3-VIB-12.7}%
\noindent\textbf{Lemma 12.7.---}
\emph{Let \(R\) be a discrete valuation ring, \(K\) its
fraction field, and \(\pi\) a uniformizer. Let \(A\) be a flat
\(R\)-algebra and let \(M\) be an \(A\)-module flat over \(R\).
Suppose that}
\begin{enumerate}
\item[\textup{(i)}]\phantomsection\label{lem:sga3-VIB-12.7-i}%
  \emph{\(M/\pi M\) is flat over
  \(\overline A=A/\pi A\),}

\item[\textup{(ii)}]\phantomsection\label{lem:sga3-VIB-12.7-ii}%
  \emph{\(M\otimes_RK\) is flat over \(A\otimes_RK\).}
\end{enumerate}
\emph{Then \(M\) is flat over \(A\).}

By the flatness criterion in the nilpotent case
\textup{(cf.\ [BAC], III \S~\SGARefLink{prop:sga3-VIB-5.2}{5.2}, Theorem~1)},
\[
  M/\pi^nM
\]
is flat over \(A/\pi^nA\) for every
\(n\in\mathbb N^*\). It then follows from \textup{[RG71]},
II, Lemma~1.4.2.1, that \(M\) is flat over \(A\). For the
reader's convenience, let us briefly indicate the proof. Put
\[
  s=\pi^n,\qquad P=M/sM.
\]
Since \(P\) is flat over \(A/sA\), and \(A/sA\) has
projective dimension \(1\) over \(A\), the spectral sequence
of composite functors shows that \(P\) has
Tor-dimension at most \(1\) over \(A\). Since \(M\) is flat
over \(R\), and thus has no \(\pi\)-torsion, \(M_K/M\) is the
direct limit of the \(A\)-modules
\[
  \pi^{-n}M/M\simeq M/\pi^nM.
\]
Consequently \(M_K/M\) also has Tor-dimension at most \(1\).
From the exact sequence
\phantomsection\label{form:sga3-VIB-12.7-MK-exact-sequence}%
\[
  0\longrightarrow M\longrightarrow M_K
    \longrightarrow M_K/M\longrightarrow0
\]
and the hypothesis that \(M_K\) is flat over \(A_K\), hence
over \(A\), it follows that \(M\) is flat.

For completeness, let us also give the following simpler
proof, pointed out by O.~Gabber. Let \(I\) be a finitely
generated ideal of \(A\). We must show that
\[
  u:M\otimes_AI\longrightarrow M
\]
is injective. By
\hyperref[lem:sga3-VIB-12.7-ii]{\textup{(ii)}},
\(u\otimes_RK\) is injective, so \(\operatorname{Ker}(u)\)
is an \(R\)-module of \(\pi\)-torsion. It therefore suffices
to show that the \(\pi\)-torsion submodule of
\(M\otimes_AI\) is zero. This submodule is a quotient of
\(\operatorname{Tor}_1^A(M,I/\pi I)\), as is seen by
tensoring with \(M\) the exact sequence
\phantomsection\label{form:sga3-VIB-12.7-I-exact-sequence}%
\[
  0\longrightarrow I\xrightarrow{\ \pi\ }I
    \longrightarrow I/\pi I\longrightarrow0.
\]
On the other hand, since \(M\) has no \(\pi\)-torsion
\textup{(being flat over \(R\))}, one has
\[
  \operatorname{Tor}_i^A(M,\overline A)=0
  \qquad(i\geq1).
\]
Consequently, if \(P_\bullet\) is a projective resolution of
the \(A\)-module \(M\), then
\[
  P_\bullet\otimes_A\overline A
\]
is a projective resolution of the \(\overline A\)-module
\(\overline M=M/\pi M\). Thus, for every
\(\overline A\)-module \(Q\),
\[
  \operatorname{Tor}_i^A(M,Q)
    =\operatorname{Tor}_i^{\overline A}(\overline M,Q),
\]
which vanishes for \(i\geq1\), since \(\overline M\) is flat
over \(\overline A\). Hence
\(\operatorname{Tor}_1^A(M,I/\pi I)=0\), proving the lemma.

\medskip
\phantomsection\label{rem:sga3-VIB-12.8}%
\noindent\textbf{Remark 12.8.---}
Let \(S\) be a regular locally noetherian scheme of dimension
\(1\), and let \(X\) be an \(S\)-scheme that is flat,
quasi-separated, and quasi-compact over \(S\). Then
\(\mathcal A(X)\) is a flat \(\mathcal O_S\)-module. Indeed,
we may suppose that \(S\) is local. Let \(s\) be its closed
point, let \(i:X_s\hookrightarrow X\) be the inclusion, and
let \(\pi\) be a uniformizer of \(R=\mathcal O(S)\). Since
\(X\) is flat over \(S\), there is an exact sequence of
sheaves
\phantomsection\label{form:sga3-VIB-12.8-sheaf-exact-sequence}%
\[
  0\longrightarrow\mathcal O_X
    \xrightarrow{\ \pi\ }\mathcal O_X
    \longrightarrow i_*(\mathcal O_{X_s})
    \longrightarrow0.
\]
Taking global sections shows that \(\mathcal A(X)\) has no
\(\pi\)-torsion as an \(R\)-module, and hence is flat.%
Moreover, the morphism
\[
  \mathcal A((X_{\mathrm{af}})_s)
    =\mathcal A(X)/\pi\mathcal A(X)
    \longrightarrow\mathcal A(X_s),
\]
induced by \(X\to X_{\mathrm{af}}\), is injective.

We can now prove the following proposition
\textup{(cf.\ \hyperref[rem:sga3-VIB-11.11.1]{Remark 11.11.1})}.

\medskip
\phantomsection\label{prop:sga3-VIB-12.9}%
\noindent\textbf{Proposition 12.9.---}
\emph{Let \(S\) be a regular locally noetherian scheme of
dimension at most \(1\), and let \(G\) be an \(S\)-group
scheme that is flat, quasi-separated, and quasi-compact over
\(S\). Then the following conditions are equivalent:}
\begin{enumerate}
\item[\textup{(i)}]\phantomsection\label{prop:sga3-VIB-12.9-i}%
  \emph{\(G\) is affine over \(S\).}

\item[\textup{(ii)}]\phantomsection\label{prop:sga3-VIB-12.9-ii}%
  \emph{\(G\) is quasi-affine over \(S\).}

\item[\textup{(iii)}]\phantomsection\label{prop:sga3-VIB-12.9-iii}%
  \emph{\(G\) acts faithfully on a quasi-affine and flat
  \(S\)-scheme \(X\).}

\item[\textup{(iv)}]\phantomsection\label{prop:sga3-VIB-12.9-iv}%
  \emph{\(G\) acts linearly and faithfully on a
  quasi-coherent module \(\mathcal E\) that is flat over
  \(S\).}

\item[\textup{(v)}]\phantomsection\label{prop:sga3-VIB-12.9-v}%
  \emph{The morphism \(\rho:G\to G_{\mathrm{af}}\) is a
  monomorphism.}
\end{enumerate}

The implication
\hyperref[prop:sga3-VIB-12.9-i]{\textup{(i)}}\(\Rightarrow\)
\hyperref[prop:sga3-VIB-12.9-ii]{\textup{(ii)}} is clear, as
is
\hyperref[prop:sga3-VIB-12.9-ii]{\textup{(ii)}}\(\Rightarrow\)
\hyperref[prop:sga3-VIB-12.9-iii]{\textup{(iii)}} \textup{(take
\(X=G\))}.

Suppose
\hyperref[prop:sga3-VIB-12.9-iii]{\textup{(iii)}} holds.
Since \(\mathcal A(G)\) and \(\mathcal A(X)\) are flat
\(\mathcal O_S\)-modules, the argument of \hyperref[prop:sga3-VIB-11.11]{11.11} shows that
\(G\) acts faithfully on the right on \(X_{\mathrm{af}}\),
and hence acts linearly and faithfully on the left on the
flat quasi-coherent \(\mathcal O_S\)-module \(\mathcal A(X)\).

If \hyperref[prop:sga3-VIB-12.9-iv]{\textup{(iv)}} holds,
\hyperref[prop:sga3-VIB-11.6.1]{11.6.1\textup{(ii)}} shows
that the monomorphism
\[
  G\longrightarrow\operatorname{Aut}_{\mathcal O_S}(\mathcal E)
\]
factors through \(G_{\mathrm{af}}\). Hence
\(G\to G_{\mathrm{af}}\) is a monomorphism.

Finally, suppose
\hyperref[prop:sga3-VIB-12.9-v]{\textup{(v)}} holds, and let
us prove that \(\rho:G\to G_{\mathrm{af}}\) is an
isomorphism. Replacing \(S\) by one of its connected
components, we may suppose that \(S\) is irreducible, with
generic point \(\eta\). Since the formation of
\(G_{\mathrm{af}}\) commutes with flat base change,
\[
  (G_{\mathrm{af}})_\eta=(G_\eta)_{\mathrm{af}}.
\]
Therefore \(G_\eta\to(G_\eta)_{\mathrm{af}}\) is a
monomorphism, hence a closed immersion by
\hyperref[lem:sga3-VIB-12.1]{12.1}, and consequently an
isomorphism because
\[
  \mathcal O((G_\eta)_{\mathrm{af}})=\mathcal O(G_\eta).
\]
If \(S=\operatorname{Spec}(\kappa(\eta))\), the proof is
complete; we may therefore suppose that \(\dim S=1\).

Let \(s\) be a closed point of \(S\). We shall show that
\[
  G_s\longrightarrow(G_{\mathrm{af}})_s
\]
is an isomorphism and that \(\rho\) is flat at every point
of \(G_s\). We may suppose that \(S\) is local, with closed
point \(s\). The base-change morphism
\[
  \rho_s:G_s\longrightarrow(G_{\mathrm{af}})_s
\]
is a monomorphism, hence a closed immersion by
\hyperref[lem:sga3-VIB-12.1]{12.1}. Thus the morphism
\[
  \mathcal O((G_{\mathrm{af}})_s)
    \longrightarrow\mathcal O(G_s)
\]
induced by \(\rho_s\) is surjective. By the preceding remark
it is also injective, and hence is an isomorphism. In
particular, \(\rho\) is surjective.

It follows from \hyperref[lem:sga3-VIB-12.7]{Lemma 12.7} that
\(\rho:G\to G_{\mathrm{af}}\) is faithfully flat. Since \(G\)
is quasi-compact over \(S\) and \(G_{\mathrm{af}}\) is
separated over \(S\), \(\rho\) is also quasi-compact
\textup{(cf.\ EGA~I, 6.6.4)}. Therefore \(\rho\) is a
quasi-compact faithfully flat monomorphism, and is thus an
isomorphism. This proves the proposition.

Finally, we prove Proposition~2.1 of Exposé~XVII,
Appendix~III. Taking \textup{[Per76]}, Corollary~4.2.5 into
account, we have replaced ``of finite type'' in the
hypotheses by ``quasi-compact and quasi-separated''.
\textup{(If \(G\) is assumed to be of finite type,
\hyperref[lem:sga3-VIB-12.1]{12.1} may be used in place of
\emph{loc.\ cit.})}

\medskip
\phantomsection\label{prop:sga3-VIB-12.10}%
\noindent\textbf{Proposition 12.10.---}
\emph{Let \(S\) be a regular locally noetherian scheme of
dimension at most \(1\), and let \(G\) be a flat,
quasi-compact, and quasi-separated \(S\)-group scheme.}
\begin{enumerate}
\item[\textup{(i)}]\phantomsection\label{prop:sga3-VIB-12.10-i}%
  \emph{The canonical morphism
  \(\rho:G\to G_{\mathrm{af}}\) is faithfully flat and
  quasi-compact. Consequently, \(G_{\mathrm{af}}\) represents
  the \textup{(fpqc)} sheaf quotient of \(G\) by
  \(N=\operatorname{Ker}(\rho)\).}

\item[\textup{(ii)}]\phantomsection\label{prop:sga3-VIB-12.10-ii}%
  \emph{If moreover \(G\) is of finite type over \(S\), then
  \(\rho\) is of finite presentation, \(G_{\mathrm{af}}\)
  represents the \textup{(fppf)} sheaf quotient of \(G\) by
  \(N\), and \(G_{\mathrm{af}}\) is of finite type over \(S\).}

\item[\textup{(iii)}]\phantomsection\label{prop:sga3-VIB-12.10-iii}%
  \emph{Suppose moreover that \(G_\eta\) is affine for every
  maximal point \(\eta\) of \(S\). Then \(N\) is an étale
  \(S\)-group and is the unit group if \(G\) is separated
  over \(S\).}
\end{enumerate}

Since \(G_{\mathrm{af}}\) is affine, and therefore separated
over \(S\), the morphism \(\rho\) is quasi-compact
\textup{(cf.\ EGA~I, 6.6.4)}, and its kernel
\[
  N=\operatorname{Ker}(\rho)
\]
is a closed subgroup of \(G\). Replacing \(S\) by one of its
connected components, we may suppose that \(S\) is
irreducible, with generic point \(\eta\).

To prove
\hyperref[prop:sga3-VIB-12.10-i]{\textup{(i)}} and
\hyperref[prop:sga3-VIB-12.10-ii]{\textup{(ii)}}, it is
enough to show that \(\rho\) is faithfully flat. Indeed,
\SGARefLink{sga3:iv:corollary:5:1:7:1}{Exposé~IV, 5.1.7.1} then shows that \(G_{\mathrm{af}}\)
represents the \textup{(fpqc)} sheaf quotient of \(G\) by
\(N\). If \(G\) is moreover of finite type over \(S\), it is
of finite presentation because \(S\) is locally noetherian;
by \hyperref[prop:sga3-VIB-9.2]{9.2\textup{(xiii)}},
\(\rho\) is then of finite presentation, as is
\(G_{\mathrm{af}}\to S\), and \(\rho\) is therefore a
covering for the \textup{(fppf)} topology.

We may thus suppose that
\[
  S=\operatorname{Spec}(R),
\]
where \(R\) is a discrete valuation ring if \(\dim S=1\), and
is the field \(\kappa(\eta)\) if \(\dim S=0\). Let \(s\) be
the closed point of \(S\). Since formation of
\(G_{\mathrm{af}}\) commutes with flat base change, the
canonical morphism \(G_\eta\to(G_\eta)_{\mathrm{af}}\) is
identified with
\[
  \rho_\eta:G_\eta\longrightarrow(G_{\mathrm{af}})_\eta.
\]
Since
\[
  \mathcal O((G_{\mathrm{af}})_\eta)
    =\mathcal O((G_\eta)_{\mathrm{af}})
    =\mathcal O(G_\eta),
\]
\(\rho_\eta\) is faithfully flat by \textup{[Per76]}, 4.2.5
\textup{(see also the addition VI\(_A\), \hyperref[cor:sga3-VIA-6.6]{6.6})}. If
\(\dim S=1\), the same argument shows that \(\rho_s\) is
faithfully flat, because the morphism
\[
  \mathcal O((G_{\mathrm{af}})_s)
    \longrightarrow\mathcal O(G_s)
\]
is injective by \hyperref[rem:sga3-VIB-12.8]{Remark 12.8}.
Lemma~\hyperref[lem:sga3-VIB-12.7]{12.7} now shows that
\(\rho\) is faithfully flat. This proves
\hyperref[prop:sga3-VIB-12.10-i]{\textup{(i)}} and
\hyperref[prop:sga3-VIB-12.10-ii]{\textup{(ii)}}. In
particular, \(N\) is flat over \(S\).

Suppose now that \(G_\eta\) is affine. Since \(\rho_\eta\)
coincides with the canonical morphism
\[
  G_\eta\longrightarrow(G_\eta)_{\mathrm{af}},
\]
its kernel \(N_\eta\) is the unit group. Let us prove that
\(N\) is étale over \(S\). Since \(N\) is flat over \(S\),
it remains to show that \(N_s\) is étale over \(\kappa(s)\)
for every point \(s\neq\eta\) of \(S\). There is nothing to
prove when \(S=\operatorname{Spec}(\kappa(\eta))\), so we may
suppose that \(S=\operatorname{Spec}(R)\), where \(R\) is a
discrete valuation ring. Let \(s\) be the closed point of
\(S\), let \(K\) be the fraction field of \(R\), and let
\(\pi\) be a uniformizer.

Let \(x\in N_s\), and let \(U_x\) be an affine open
neighborhood of \(x\) in \(N\). Since \(N\) is flat over
\(S\), the intersection \(U_x\cap N_\eta\) is nonempty, and
hence equals \(\{\epsilon(\eta)\}\), where \(\epsilon\)
denotes the unit section. Thus
\[
  A=\mathcal O(U_x)
\]
is a flat \(R\)-algebra such that \(A_K=K\) and
\(\pi^{-1}\notin A\), since \(\pi\) belongs to the maximal
ideal of \(\mathcal O_{U_x,x}\). It follows that \(A=R\), and
hence the projection \(U_x\to S\) is an isomorphism. Thus
\(N\) is étale over \(S\). If \(G\) is also separated over
\(S\), the inverse isomorphism \(S\to U_x\) is the unit
section, since the two coincide on the dense open subset
\(\{\eta\}\) of \(S=\operatorname{Spec}(R)\). Hence \(N\) is
the unit group. This proves the proposition.

In particular, we obtain the following corollary, two other
proofs of which may be found in \textup{[An73]},
\SGARefLink{prop:sga3-VIA-2.3.1}{Proposition~2.3.1}, and \textup{[PY06]}, \SGARefLink{def:sga3-VIB-3.1}{Proposition~3.1}.

\medskip
\phantomsection\label{cor:sga3-VIB-12.10.1}%
\noindent\textbf{Corollary 12.10.1.---}
\emph{Let \(R\) be a discrete valuation ring, \(K\) its
fraction field, and \(G\) a separated, flat \(R\)-group
scheme of finite type over \(R\). If \(G_K\) is affine, then
\(G\) is affine.}

\medskip
\phantomsection\label{rem:sga3-VIB-12.10.2}%
\noindent\textbf{Remarks 12.10.2.---}
\begin{enumerate}
\item[\textup{(a)}]\phantomsection\label{rem:sga3-VIB-12.10.2-a}%
  O.~Gabber has pointed out to us examples in which \(G\) is
  a flat group of finite type over a discrete valuation ring,
  its generic fiber is an abelian variety, and the kernel
  \(N\) of \(G\to G_{\mathrm{af}}\) is not smooth.

\item[\textup{(b)}]\phantomsection\label{rem:sga3-VIB-12.10.2-b}%
  M.~Raynaud, on the other hand, gave an example, with \(S\)
  the affine plane of dimension \(2\) over a field \(k\), of
  a smooth quasi-affine \(S\)-group scheme with affine,
  connected fibers that is not affine over \(S\); see
  \textup{[Ray70a]}, \S~VII.3, p.~116.
\end{enumerate}
